\documentclass[11pt]{article}
\usepackage[margin=1in]{geometry}
\usepackage{amsmath,amssymb,amsthm,mathtools,booktabs}
\usepackage[hidelinks]{hyperref}

\newtheorem{theorem}{Theorem}[section]
\newtheorem{lemma}[theorem]{Lemma}
\newtheorem{proposition}[theorem]{Proposition}
\newtheorem{corollary}[theorem]{Corollary}
\theoremstyle{definition}
\newtheorem{definition}[theorem]{Definition}

\newcommand{\SYT}{\operatorname{SYT}}
\newcommand{\ellp}{\ell}

\title{Exact Enumeration and Fixed-Rank Asymptotics\\
of Ballot-Admissible Fibonacci Ribbon Tableaux}
\author{Alex Chengyu Li}
\date{29 August 2026}
\hypersetup{
  pdftitle={Exact Enumeration and Fixed-Rank Asymptotics of Ballot-Admissible Fibonacci Ribbon Tableaux},
  pdfauthor={Alex Chengyu Li}
}

\begin{document}
\maketitle

\begin{center}
\begin{minipage}{0.88\textwidth}
\textbf{Formalization status.}
All 47 numbered theorem, corollary, lemma, and equation labels have
premise-free publication endpoints in the accompanying
\href{https://github.com/crabsatellite/fibonacci-ribbon-ballot-enumeration}{Lean~4 development}.
The complete publication root is checked with \texttt{--trust=0}; its axiom
audit reports only \texttt{propext}, \texttt{Classical.choice}, and
\texttt{Quot.sound}.
\end{minipage}
\end{center}
\medskip

\begin{abstract}
Tenn asked for the enumeration of ballot-admissible Fibonacci ribbon
tableaux beyond the three-letter case.  We solve this problem and
simultaneously cover every alphabet size $n\ge2$ by proving the following
exact formula.  If $b_{n,k}$ denotes the number of ballot-admissible
$\operatorname{Fib}(n,k)$ ribbons and $f^\lambda$ is the number of standard
Young tableaux of shape $\lambda$, then
\[
 b_{n,k}=\sum_{j=0}^{\lfloor k/2\rfloor}(-1)^j
 \binom{k-j}{j}
 \sum_{\substack{\lambda\vdash k-2j\\\ellp(\lambda)\le n}}f^\lambda.
\]
The proof encodes the columns as a dominant type-$A$ walk alternating between
the defining representation and its dual.  The unrestricted mixed walk has
the bounded-height standard-tableau count, while each forbidden adjacency is
a ballot-neutral block; endpoint-preserving contraction and
inclusion--exclusion give both the formula and a highest-weight refinement in
explicit Schur coefficients.

The resulting generating-function substitution, together with Gessel's
finite-dimensional Bessel system and Regev's strip asymptotics, yields an
explicit constant $K_n>0$ such that, for every fixed $n\ge3$,
\[
 b_{n,k}\sim K_n
 \left(\frac{n+\sqrt{n^2-4}}2\right)^k k^{-n(n-1)/4}.
\]
For $n\ge k$, the count stabilizes to involutions with no adjacent
transposition.  Their number is asymptotic to $e^{-1}$ times the involution
number, proving a conjecture recorded in OEIS~A170941.  More generally, for
each fixed $r$, the defect on the line $n=k-r$ is eventually a polynomial in
$k$ of degree $r-1$ with leading coefficient $I_{r-1}/(r-1)!$, where $I_m$
is the $m$th involution number.
\end{abstract}

\section{Introduction}

Fibonacci ribbons are alternating-height skew tableaux introduced by
Donnelly, Dunkum, Li\v{s}kov\'a, and Nance in their study of symmetric
Fibonaccian distributive lattices and special-linear representations
\cite{DonnellyEtAl2023}.  A ribbon is ballot-admissible when it is a
Littlewood--Richardson tableau.  For three letters, empirical data suggested
the peakless Motzkin numbers.  Tenn proved that conjecture and constructed an
explicit bijection, then left the enumeration of
$\operatorname{Fib}(n,k)$ for $n\ge4$ as an open problem
\cite{Tenn2026}.

The main result of this paper is an exact finite formula for every $n$ and
$k$.  The derivation uses two structural facts.  First, if the
single local semistandard obstruction is temporarily allowed, the alternating
short and tall columns form a mixed defining/dual branching walk.  Although
the endpoint distributions depend on which factors are dual, their total
multiplicity does not.  Second, each forbidden adjacent pair reads exactly
$12\cdots n$, so it may be deleted without changing any ballot difference.
The obstruction sets are matchings of a path, and inclusion--exclusion gives
the answer.

The contraction preserves the final dominant weight, which yields a
highest-weight refinement in explicit Schur coefficients.  The ordinary
generating function is obtained from the bounded-tableau series by a rational
substitution.  Combining this substitution with Gessel's finite-dimensional
Bessel system and Regev's strip asymptotics gives, for every fixed $n\ge3$,
an explicit constant $K_n$ such that
\[
 b_{n,k}\sim K_n
 \left(\frac{n+\sqrt{n^2-4}}2\right)^k
 k^{-n(n-1)/4}.
\]

The stable range leads to a second theme.  The sum of the numbers of standard
tableaux of all shapes of size $k$ is the involution number.  The same
path-matching inclusion--exclusion excludes the cycles $(i,i+1)$.  This
identifies the stable ribbon numbers with OEIS~A170941 and supplies a short
Poisson-limit proof of the conjectural limiting proportion recorded there
\cite{OEISA170941}.

The approach to stability is also explicit.  On every fixed-distance line
$n=k-r$, the deficit from the stable count is a finite sum of bounded-tail
tableau numbers and is eventually a polynomial in $k$ of degree $r-1$.
Together with the fixed-rank asymptotic, this also shows that the two natural
iterated density limits do not commute.

\section{Fibonacci ribbons and column parameters}

Fix integers $n\ge2$ and $k\ge0$.  The Fibonacci ribbon shape has $k$
columns.  From right to left their heights alternate between $1$ and $n-1$,
starting with height $1$.  Entries lie in $[n]=\{1,\ldots,n\}$, columns are
strictly increasing, and rows are weakly increasing.  The linear word is read
column by column from right to left and within a column from top to bottom.
The tableau is ballot-admissible if every prefix $w'$ of this word satisfies
\[
 \#_1(w')\ge\#_2(w')\ge\cdots\ge\#_n(w').
\]
These are the conventions of \cite{DonnellyEtAl2023,Tenn2026}.

Parameterize the $c$th column from the right by $a_c\in[n]$.  If $c$ is odd,
the column is the singleton $a_c$.  If $c$ is even, it contains every element
of $[n]$ except
\[
 y_c=n+1-a_c,
\]
in increasing order.  The Fibonaccian-string condition from
\cite[Theorem~4.2]{DonnellyEtAl2023}, translated from left-to-right to
right-to-left order, becomes
\begin{equation}\label{eq:local-obstruction}
 (a_c,a_{c+1})\ne(1,n).
\end{equation}
Thus the local obstruction is unique.

For a prefix $w$, define its difference vector
\[
 d(w)=(d_1,\ldots,d_{n-1}),\qquad
 d_i=\#_i(w)-\#_{i+1}(w).
\]
The ballot condition is $d(w)\in\mathbb Z_{\ge0}^{n-1}$ at every prefix.
Put
\[
 v_1=e_1,\qquad
 v_x=-e_{x-1}+e_x\ (1<x<n),\qquad
 v_n=-e_{n-1}.
\]

\begin{lemma}[Literal column walk]\label{lem:column-walk}
Reading a singleton $x$ changes $d$ by $v_x$.  Reading a tall column omitting
$y$ changes $d$ by $-v_y$.  In either case the column may be appended to a
ballot prefix exactly when the resulting difference vector is nonnegative.
\end{lemma}

\begin{proof}
The singleton statement follows directly from the definition of the
differences.  A full increasing block $12\cdots n$ returns every difference
to its initial value and stays nonnegative when started from a nonnegative
vector.  Omitting $y$ therefore has net increment $-v_y$.

It remains to check internal prefixes.  Before the omitted position, the
increasing letters merely move one unit of surplus to the right.  If $y<n$,
the only possible new violation occurs when $y+1$ is read, and it occurs
exactly when the final coordinate required by $d-v_y\ge0$ is absent.  If
$y=n$, there is no later letter and no additional condition.  Hence endpoint
nonnegativity is necessary and sufficient.
\end{proof}

Let $b_{n,k}$ denote the desired ribbon count.  Let $u_{n,k}$ denote the
number obtained from the same ballot column words when
Equation~\eqref{eq:local-obstruction} is not imposed.

\section{The unrestricted mixed branching walk}

Let $D_n=\mathbb Z_{\ge0}^{n-1}$.  On functions on $D_n$, let $A$ be the
adjacency operator for the valid steps $v_x$, and let $B$ be the adjacency
operator for the valid steps $-v_y$.  All sums below are finite at every
state.

\begin{lemma}[Mixed branching]\label{lem:mixed-branching}
The operators $A$ and $B$ commute, and
\[
 A\mathbf1=B\mathbf1,
\]
where $\mathbf1$ is the constant-one function on $D_n$.
\end{lemma}

\begin{proof}
The vectors $v_x-v_y$ are distinct when $x\ne y$.  For such a pair, the
validity conditions of the two steps are independent except when $x=y+1$.
In that exceptional case, both orders are valid exactly when $d_y\ge2$:
each step lowers the same boundary coordinate by one.  Hence a step $v_x$
followed by $-v_y$ is valid if and only if the reverse order is valid for
every $x\ne y$, so the off-diagonal matrix entries of $AB$ and $BA$ agree.

For a zero-displacement two-step walk, the two letters must have $x=y$.
The number of valid $A$-moves from $d$ is
\[
 1+\#\{i:d_i>0\};
\]
the always-valid move is $x=1$.  The number of valid $B$-moves is the same,
with always-valid move $y=n$.  Hence the diagonal entries also agree and
$AB=BA$.  The same outdegree calculation gives
$A\mathbf1=B\mathbf1$.
\end{proof}

\begin{theorem}[Unrestricted count]\label{thm:unrestricted}
For every $n\ge2$ and $k\ge0$,
\begin{equation}\label{eq:unrestricted}
 u_{n,k}=\sum_{\substack{\lambda\vdash k\\\ellp(\lambda)\le n}}f^\lambda,
\end{equation}
where $f^\lambda=|\SYT(\lambda)|$.
\end{theorem}

\begin{proof}
There are $a=\lceil k/2\rceil$ defining steps and
$c=\lfloor k/2\rfloor$ dual steps.  If $\delta_0$ is evaluation at the zero
weight, the total unrestricted count is
\[
 \delta_0 A^aB^c\mathbf1.
\]
Lemma~\ref{lem:mixed-branching} gives, inductively,
$B^r\mathbf1=A^r\mathbf1$: indeed
\[
 B^{r+1}\mathbf1=BA^r\mathbf1=A^rB\mathbf1=A^{r+1}\mathbf1.
\]
Consequently $u_{n,k}=\delta_0A^k\mathbf1$.

A path of $k$ defining steps from zero is a word
$x_1\cdots x_k$ on $[n]$ whose prefix multiplicities are weakly decreasing.
Place $j$ in row $x_j$.  Rows increase because positions are inserted in
order, and the prefix inequalities say precisely that every entry has a cell
above it when required.  This is a standard Young tableau with at most $n$
rows.  Reading the row containing each successive entry reverses the
construction.  Summing over shapes proves Equation~\eqref{eq:unrestricted}.
\end{proof}

By the hook-length formula,
\begin{equation}\label{eq:hook}
 f^\lambda=\frac{k!}{\prod_{z\in\lambda}h(z)}
 \qquad(\lambda\vdash k),
\end{equation}
so Equation~\eqref{eq:unrestricted} is directly computable.
Robinson--Schensted also identifies $u_{n,k}$ with the number of involutions
of $[k]$ whose longest decreasing subsequence has length at most $n$.

\section{Exact enumeration by neutral-block contraction}

\begin{lemma}[Neutral bad pairs]\label{lem:neutral-pairs}
Suppose a specified adjacency satisfies $(a_c,a_{c+1})=(1,n)$.  Its two
columns have combined linear word $12\cdots n$.  Deleting the two columns is
a bijection between unrestricted ballot column words having that specified
bad pair and unrestricted ballot column words of length $k-2$.

More generally, if $j$ pairwise disjoint adjacencies are specified to be bad,
simultaneously deleting them is a bijection with unrestricted objects of
length $k-2j$.
\end{lemma}

\begin{proof}
If $c$ is odd, the two blocks are the singleton $1$ and the tall column
$2\cdots n$.  If $c$ is even, they are $1\cdots(n-1)$ and the singleton $n$.
Their concatenation is $12\cdots n$ in both cases.  This block returns every
difference to its initial value, and every internal prefix is ballot.
Deletion and insertion are inverse operations.  Removing two columns
preserves the parity of every later column, so the argument applies
simultaneously to any disjoint collection.
\end{proof}

Bad adjacencies cannot overlap, since a shared parameter would have to equal
both $1$ and $n$.  Sets of $j$ possible bad locations are therefore the
$j$-edge matchings of a path on $k$ vertices, of which there are
\begin{equation}\label{eq:path-matchings}
 \binom{k-j}{j}.
\end{equation}

\begin{theorem}[Fibonacci-ribbon enumeration]\label{thm:main}
For every $n\ge2$ and $k\ge0$,
\begin{equation}\label{eq:main}
 \boxed{
 b_{n,k}=\sum_{j=0}^{\lfloor k/2\rfloor}(-1)^j
 \binom{k-j}{j}
 \sum_{\substack{\lambda\vdash k-2j\\\ellp(\lambda)\le n}}f^\lambda.}
\end{equation}
Equations~\eqref{eq:hook} and~\eqref{eq:main} give a finite explicit formula.
\end{theorem}

\begin{proof}
Apply inclusion--exclusion to the $k-1$ bad-adjacency events.  An intersection
is empty unless its locations form a matching.  For a matching of size $j$,
Lemma~\ref{lem:neutral-pairs} identifies the intersection with the
unrestricted objects counted by $u_{n,k-2j}$.  Use
Equations~\eqref{eq:path-matchings} and~\eqref{eq:unrestricted}.
\end{proof}

The same contraction retains the final dominant weight and therefore admits
a highest-weight refinement.  For a dominant difference vector
$d=(d_1,\ldots,d_{n-1})$, let
\[
 \mu_i=d_i+d_{i+1}+\cdots+d_{n-1}\quad(1\le i<n),
 \qquad \mu_n=0.
\]
Thus $\mu$ is the normalized partition representing the corresponding
$\mathfrak{sl}_n$ highest weight.  Let $b_{n,k}^{\mu}$ count admissible
ribbons whose column walk ends at this weight.  Put
\[
 a=\left\lceil\frac{k}{2}\right\rceil,
 \qquad c=\left\lfloor\frac{k}{2}\right\rfloor.
\]
For nonnegative $r,s$, define
\begin{equation}\label{eq:mixed-schur-multiplicity}
 M_{r,s}^{(n)}(\mu)=
 \begin{cases}
 \left\langle p_1^r e_{n-1}^s,
 s_{\mu+(q^n)}\right\rangle,
 &q=\dfrac{r+(n-1)s-|\mu|}{n}\in\mathbb Z_{\ge0},\\[6pt]
 0,&\text{otherwise},
 \end{cases}
\end{equation}
where the bracket denotes the Schur coefficient, and
$\mu+(q^n)=(\mu_1+q,\ldots,\mu_n+q)$.

\begin{theorem}[Highest-weight refinement]\label{thm:highest-weight}
For every normalized dominant partition $\mu$,
\begin{equation}\label{eq:highest-weight}
 \boxed{
 b_{n,k}^{\mu}=
 \sum_{j=0}^{\lfloor k/2\rfloor}(-1)^j\binom{k-j}{j}
 M_{a-j,c-j}^{(n)}(\mu).}
\end{equation}
Summing Equation~\eqref{eq:highest-weight} over $\mu$ recovers
Theorem~\ref{thm:main}.
\end{theorem}

\begin{proof}
Let $V$ be the defining $SL_n$-module.  As an $SL_n$-module,
$V^*\cong\bigwedge^{n-1}V$.  Hence the polynomial $GL_n$ character whose
restriction describes $V^{\otimes r}\otimes(V^*)^{\otimes s}$ is
$p_1^r e_{n-1}^s$.  A polynomial $GL_n$ constituent $S_\lambda V$ restricts
to normalized $SL_n$ highest weight $\mu$ exactly when
$\lambda=\mu+(q^n)$ for some $q\ge0$.  Homogeneous degree forces
$|\lambda|=r+(n-1)s$, so $q$ is the unique integer in
Equation~\eqref{eq:mixed-schur-multiplicity}.  The Pieri branching paths to
$\mu$ are therefore counted by $M_{r,s}^{(n)}(\mu)$
\cite{StanleyEC2}.

Every selected bad adjacency contains one defining and one dual column.
Deleting its full-alphabet block changes $(a,c)$ to $(a-1,c-1)$ but has zero
difference displacement, so it preserves the final highest weight.  The
simultaneous contraction of $j$ disjoint bad pairs therefore has endpoint
multiplicity $M_{a-j,c-j}^{(n)}(\mu)$.  Endpointwise inclusion--exclusion
proves Equation~\eqref{eq:highest-weight}.  Summation over all endpoints is
the unrestricted endpoint functional used in Theorem~\ref{thm:unrestricted},
which gives the last assertion.
\end{proof}

The first rows, beginning at $k=0$, are
\begin{center}
\small
\begin{tabular}{c|rrrrrrrrrrrrr}
\toprule
$n\backslash k$&0&1&2&3&4&5&6&7&8&9&10&11&12\\
\midrule
2&1&1&1&1&1&1&1&1&1&1&1&1&1\\
3&1&1&1&2&4&8&17&37&82&185&423&978&2283\\
4&1&1&1&2&5&12&31&82&229&646&1891&5574&16883\\
5&1&1&1&2&5&13&36&105&321&1018&3334&11216&38635\\
6&1&1&1&2&5&13&37&111&355&1179&4099&14665&54299\\
\bottomrule
\end{tabular}
\end{center}
For $n=3$, Theorem~\ref{thm:main} recovers the peakless Motzkin numbers and
therefore agrees with the bijection of \cite{Tenn2026}.

\section{Generating functions and finite-rank structure}

Put
\[
 U_n(z)=\sum_{k\ge0}u_{n,k}z^k,
 \qquad B_n(z)=\sum_{k\ge0}b_{n,k}z^k.
\]
For fixed $n$ these are convergent near the origin.  They may also be read as
formal power series.

\begin{theorem}[Exact substitution]\label{thm:substitution}
For every $n\ge2$,
\begin{equation}\label{eq:substitution}
 B_n(z)=\frac{1}{1+z^2}
 U_n\!\left(\frac{z}{1+z^2}\right).
\end{equation}
Consequently $B_n(z)$ is D-finite for every fixed $n$.
\end{theorem}

\begin{proof}
In Theorem~\ref{thm:main}, put $m=k-2j$ and sum over $k$.  The elementary
identity
\[
 \sum_{j\ge0}\binom{m+j}{j}t^j=(1-t)^{-m-1}
\]
gives
\begin{align*}
 B_n(z)
 &=\sum_{m\ge0}u_{n,m}z^m
   \sum_{j\ge0}(-1)^j\binom{m+j}{j}z^{2j}\\
 &=\frac{1}{1+z^2}
   \sum_{m\ge0}u_{n,m}
   \left(\frac{z}{1+z^2}\right)^m.
\end{align*}
For fixed $n$, the bounded-height tableau sequence $u_{n,m}$ is
$P$-recursive; equivalently its generating function is D-finite
\cite{Gessel1990}.  D-finite series are closed under multiplication by a
rational function and algebraic substitution, proving the last assertion.
\end{proof}

The substitution is invertible as a formal series.  If
\[
 \psi(w)=\frac{1-\sqrt{1-4w^2}}{2w}=w+w^3+2w^5+\cdots,
\]
then $\psi/(1+\psi^2)=w$ and
\begin{equation}\label{eq:inverse-substitution}
 U_n(w)=(1+\psi(w)^2)B_n(\psi(w)).
\end{equation}
Thus the ribbon and bounded-involution enumerations determine one another.

For $n=3$, $u_{3,k}$ is the Motzkin number, so
$U_3=1+zU_3+z^2U_3^2$.  Equation~\eqref{eq:substitution} simplifies to
\[
 B_3(z)=1+zB_3(z)+z^2B_3(z)(B_3(z)-1),
\]
which is exactly the peakless-Motzkin decomposition used in
\cite{Tenn2026}.  For $n=2$, the same substitution reduces to
$B_2(z)=1/(1-z)$, agreeing with $b_{2,k}=1$.

The substitution in fact transfers the complete leading strip asymptotic in
every fixed rank.  Put
\[
 \beta_n=\frac{n(n-1)}4,
 \qquad
 \alpha_n=\frac{n+\sqrt{n^2-4}}2,
\]
and
\begin{equation}\label{eq:regev-constant}
 C_n=\frac{n^{\beta_n}}{n!}
 \prod_{j=1}^n\frac{\Gamma(1+j/2)}{\Gamma(3/2)}.
\end{equation}

\begin{theorem}[Fixed-rank asymptotic]\label{thm:fixed-rank-asymptotic}
For every fixed $n\ge3$,
\begin{equation}\label{eq:fixed-rank-asymptotic}
 b_{n,k}\sim K_n\alpha_n^k k^{-\beta_n},
 \qquad
 K_n=C_n\frac{\alpha_n}{n}
 \left(\frac{\sqrt{n^2-4}}{n}\right)^{\beta_n-1}.
\end{equation}
For $n=2$, one has $b_{2,k}=1$ exactly.
\end{theorem}

\begin{proof}
Regev's strip theorem, with the Mehta integral evaluated explicitly, gives
\begin{equation}\label{eq:regev-leading}
 u_{n,k}\sim C_n n^k k^{-\beta_n}
\end{equation}
\cite{Regev1981,Matsumoto2008}.  It remains to justify that no other
singularity becomes dominant under Equation~\eqref{eq:substitution}.

Let
\[
 Y_n(x)=\sum_{k\ge0}u_{n,k}\frac{x^k}{k!},
 \qquad
 J_r(x)=\sum_{m\ge0}\frac{x^{2m+r}}{m!(m+r)!}.
\]
Write $d=\lfloor n/2\rfloor$.  Gessel's determinant, in the simplified form
of Bergeron and Gascon, places $Y_n$ in the $\mathbb C(x)$-span
\[
 \{J_0^aJ_1^{d-a}:0\le a\le d\}
\]
when $n=2d$, and in $e^x$ times this span when $n=2d+1$
\cite{Gessel1990,BergeronGascon2000}.  The coefficients may be taken to be
Laurent polynomials in $x$: the recurrence reducing every $J_r$ to $J_0,J_1$
introduces only powers of $x^{-1}$.  The identities
\[
 J_0'=2J_1,
 \qquad
 J_1'=2J_0-x^{-1}J_1
\]
therefore give a basis vector $\mathbf F$ satisfying an exact system
\begin{equation}\label{eq:bessel-system}
 \mathbf F'(x)=(M_0+x^{-1}M_1)\mathbf F(x)
\end{equation}
for constant matrices $M_0,M_1$.  In the basis generated by
$J_0+J_1$ and $J_0-J_1$, the eigenvalues of $M_0$ are
\begin{equation}\label{eq:bessel-scales}
 n,n-4,n-8,\ldots,n-4d.
\end{equation}
Indeed the two generators have limiting derivative eigenvalues $2$ and
$-2$, and the odd case adds the derivative eigenvalue $1$ of $e^x$.

Expand $\mathbf F=\sum_{k\ge0}\mathbf f_kx^k/k!$ and put
$\widehat{\mathbf F}(z)=\sum_{k\ge0}\mathbf f_kz^k$.  Multiplying
Equation~\eqref{eq:bessel-system} by $x$ and comparing coefficients gives
\[
 (kI-M_1)\mathbf f_k=kM_0\mathbf f_{k-1}.
\]
After summation,
\begin{equation}\label{eq:ordinary-bessel-system}
 z(I-zM_0)\widehat{\mathbf F}'(z)
 =(M_1+zM_0)\widehat{\mathbf F}(z)-M_1\mathbf f_0.
\end{equation}
Multiplication by a Laurent monomial in the exponential-generating variable
becomes a finite combination of shifts, derivatives, and integrals of the
ordinary generating function; it introduces no new nonzero finite
singularity.  Consequently every nonzero finite singularity of $U_n$ belongs to
\[
 \left\{\frac1{n-4j}:0\le j\le d,\ n-4j\ne0\right\}.
\]
These are regular singularities, so their coefficient contributions are
finite sums of the standard power--logarithm forms.
Equation~\eqref{eq:regev-leading} and Pringsheim's theorem identify $1/n$ as
the leading positive singularity and fix its leading coefficient term.
When $n$ is even, $-1/n$ is the only other point of the same modulus; its
coefficient contribution is of lower order, since otherwise the even and
odd subsequences would contradict the common asymptotic
Equation~\eqref{eq:regev-leading}.  Every remaining scale has modulus
strictly smaller than $n$.

Let $w(z)=z/(1+z^2)$ and $\rho_n=\alpha_n^{-1}$.  The positive singularity
$1/n$ has the simple preimage $\rho_n$.  For a nonzero scale $\lambda$ with
$|\lambda|<n$, the closest preimage of $1/\lambda$ has modulus
\[
 \frac{2}{|\lambda|+\sqrt{\lambda^2-4}}>\rho_n
\]
when $|\lambda|\ge2$, and modulus one when $|\lambda|\le2$.  Thus every
smaller scale has all its preimages at strictly larger modulus.  The possible
point $-1/n$ maps to $-\rho_n$ and remains of lower coefficient order.  At
the positive point,
\[
 \frac{\rho_n w'(\rho_n)}{1/n}=\frac{\sqrt{n^2-4}}n,
 \qquad
 \frac1{1+\rho_n^2}=\frac{\alpha_n}{n}.
\]
The standard power--logarithm transfer theorem
\cite[Chapter~VI]{FlajoletSedgewick2009}, including its logarithmic form when
$\beta_n$ is an integer, therefore multiplies the coefficient constant $C_n$
by the two factors displayed in
Equation~\eqref{eq:fixed-rank-asymptotic}.  This proves the result.
\end{proof}

The first dimension left open by Tenn admits a particularly explicit form.
Let $C_r=\frac1{r+1}\binom{2r}{r}$.

\begin{corollary}[The four-letter case]\label{cor:n-four}
Define
\[
 c_{2r}=C_rC_{r+1},\qquad c_{2r+1}=C_{r+1}^2.
\]
Then
\begin{equation}\label{eq:n-four-sum}
 b_{4,k}=\sum_{j=0}^{\lfloor k/2\rfloor}
 (-1)^j\binom{k-j}{j}c_{k-2j}.
\end{equation}
Moreover,
\begin{equation}\label{eq:n-four-asymptotic}
 b_{4,k}\sim
 \frac{6(2+\sqrt3)}{\pi}
 (2+\sqrt3)^k k^{-3}.
\end{equation}
\end{corollary}

\begin{proof}
The bounded-height tableau formula
\[
 u_{4,2r}=C_rC_{r+1},\qquad u_{4,2r+1}=C_{r+1}^2
\]
is classical \cite{GouyouBeauchamps1989}.  Substitution in
Theorem~\ref{thm:main} gives Equation~\eqref{eq:n-four-sum}.
Equation~\eqref{eq:regev-constant} gives $C_4=32/\pi$.  Applying
Theorem~\ref{thm:fixed-rank-asymptotic}, with
$\beta_4=3$, $\alpha_4=2+\sqrt3$, and
$\sqrt{4^2-4}/4=\sqrt3/2$, gives
Equation~\eqref{eq:n-four-asymptotic}.
\end{proof}

The next rank also admits a closed recurrence and a complete transferred
asymptotic.  Put
\[
 \alpha_5=\frac{5+\sqrt{21}}2.
\]

\begin{corollary}[The five-letter case]\label{cor:n-five}
Let $c_m=u_{5,m}$, with $c_0=c_1=1$, $c_2=2$.  For $m\ge3$ these numbers
satisfy
\begin{align}\label{eq:height-five-recurrence}
 &(m+4)(m+6)c_m-(3m^2+17m+15)c_{m-1}\\
 &\quad{}-(m-1)(13m+9)c_{m-2}
 +15(m-1)(m-2)c_{m-3}=0.\notag
\end{align}
The ribbon number is
\begin{equation}\label{eq:n-five-sum}
 b_{5,k}=\sum_{j=0}^{\lfloor k/2\rfloor}
 (-1)^j\binom{k-j}{j}c_{k-2j},
\end{equation}
and
\begin{equation}\label{eq:n-five-asymptotic}
 b_{5,k}\sim
 \frac{1323\alpha_5}{8\pi}\,
 \alpha_5^k k^{-5}.
\end{equation}
\end{corollary}

\begin{proof}
The height-five recurrence is the bounded-tableau recurrence of
Bergeron and Gascon \cite{BergeronGascon2000};
Equation~\eqref{eq:n-five-sum} is Theorem~\ref{thm:main}.
Equation~\eqref{eq:regev-constant} gives $C_5=9375/(8\pi)$.  Now
Theorem~\ref{thm:fixed-rank-asymptotic}, with $\beta_5=5$ and local scale
$\sqrt{21}/5$, gives
\[
 \frac{9375}{8\pi}\frac{\alpha_5}{5}
 \left(\frac{\sqrt{21}}5\right)^4
 =\frac{1323\alpha_5}{8\pi},
\]
which is Equation~\eqref{eq:n-five-asymptotic}.
\end{proof}

The next rank provides another explicit recurrence.  Put
$\alpha_6=3+2\sqrt2$.

\begin{corollary}[The six-letter case]\label{cor:n-six}
Let $c_m=u_{6,m}$, with $c_0=c_1=1$, $c_2=2$, and $c_3=4$.  For $m\ge4$,
\begin{align}\label{eq:height-six-recurrence}
 &(m+5)(m+8)(m+9)c_m
 -4(5m^2+46m+84)c_{m-1}\notag\\
 &\quad{}-4(m-1)(10m^2+58m+33)c_{m-2}
 +144(m-1)(m-2)c_{m-3}\notag\\
 &\quad{}+144(m-1)(m-2)(m-3)c_{m-4}=0.
\end{align}
Moreover,
\begin{equation}\label{eq:n-six-sum}
 b_{6,k}=\sum_{j=0}^{\lfloor k/2\rfloor}
 (-1)^j\binom{k-j}{j}c_{k-2j},
\end{equation}
and
\begin{equation}\label{eq:n-six-asymptotic}
 b_{6,k}\sim
 \frac{3\,2^{57/4}\alpha_6}{\pi^{3/2}}
 \alpha_6^k k^{-15/2}.
\end{equation}
\end{corollary}

\begin{proof}
The differential equation for $Y_6$ in Bergeron and Gascon
\cite{BergeronGascon2000}, after coefficient extraction, is precisely
Equation~\eqref{eq:height-six-recurrence}.  Equation~\eqref{eq:n-six-sum}
is Theorem~\ref{thm:main}.  Finally,
Equation~\eqref{eq:regev-constant} gives
\[
 C_6=\frac{3}{4}\frac{6^{15/2}}{\pi^{3/2}}.
\]
Theorem~\ref{thm:fixed-rank-asymptotic} has
$\beta_6=15/2$ and local scale $\sqrt{32}/6=2\sqrt2/3$; simplifying
\[
 C_6\frac{\alpha_6}{6}
 \left(\frac{2\sqrt2}{3}\right)^{13/2}
\]
gives the constant in Equation~\eqref{eq:n-six-asymptotic}.
\end{proof}

\section{Stable involutions}

Let $I_k$ be the number of involutions of $[k]$.  Robinson--Schensted gives a
bijection between involutions and standard Young tableaux, because an
involution has equal insertion and recording tableaux
\cite{Sagan2001,StanleyEC2}.  Thus
\begin{equation}\label{eq:involution-tableaux}
 I_k=\sum_{\lambda\vdash k}f^\lambda,
 \qquad I_k=I_{k-1}+(k-1)I_{k-2}.
\end{equation}

\begin{corollary}[Stable range]\label{cor:stable}
If $n\ge k$, then $b_{n,k}$ is the number $a_k$ of involutions of $[k]$
having no cycle $(i,i+1)$.  Equivalently, it is the number of matchings in
the complement of the path $P_k$.  In particular,
\begin{equation}\label{eq:stable-ie}
 a_k=\sum_{j=0}^{\lfloor k/2\rfloor}(-1)^j
 \binom{k-j}{j}I_{k-2j}.
\end{equation}
The initial values are
\[
1,1,1,2,5,13,37,112,363,1235,4427,16526,\ldots.
\]
\end{corollary}

\begin{proof}
When $n\ge k$, every partition of $k-2j$ has at most $n$ rows, so
Theorem~\ref{thm:main} and Equation~\eqref{eq:involution-tableaux} give
Equation~\eqref{eq:stable-ie}.  To exclude adjacent transpositions directly,
apply inclusion--exclusion to the path edges $(i,i+1)$.  A nonempty
intersection requires a matching of $j$ path edges.  Once those $j$
transpositions are fixed, the remaining $k-2j$ labels support an arbitrary
involution.  There are $\binom{k-j}{j}$ such matchings, giving the same sum.
\end{proof}

The first finite-rank corrections below stability are also explicit.  For
$s,m\ge0$, put
\begin{equation}\label{eq:tail-tableaux}
 E_s(m)=\sum_{\substack{\lambda\vdash m\\m-\lambda_1\le s}}f^\lambda.
\end{equation}

\begin{corollary}[Near-stable defect]\label{cor:near-stable}
Fix $r\ge1$ and let $k\ge r+2$.  Then
\begin{equation}\label{eq:near-stable-defect}
 a_k-b_{k-r,k}
 =\sum_{0\le j<r/2}(-1)^j\binom{k-j}{j}
 E_{r-2j-1}(k-2j).
\end{equation}
For each fixed $r$, this defect is eventually a polynomial in $k$ of degree
$r-1$ with leading coefficient $I_{r-1}/(r-1)!$.  The first four strips are
\begin{align}\label{eq:first-near-stable-strips}
 b_{k-1,k}&=a_k-1,\notag\\
 b_{k-2,k}&=a_k-k,\notag\\
 b_{k-3,k}&=a_k-(k-1)(k-2),\\
 b_{k-4,k}&=a_k-\frac{(k-2)(2k^2-8k+3)}3.\notag
\end{align}
\end{corollary}

\begin{proof}
Subtract Theorem~\ref{thm:main}, with $n=k-r$, from the stable formula
Equation~\eqref{eq:stable-ie}.  At the term $m=k-2j$, the missing tableau
shapes have more than $k-r$ rows.  Conjugating their diagrams preserves
$f^\lambda$ and turns this condition into
\[
 \lambda_1>k-r=m-(r-2j),
\]
so their total is exactly $E_{r-2j-1}(m)$.  There is no defect once
$2j\ge r$, proving Equation~\eqref{eq:near-stable-defect}.

For fixed $s$, write every shape in Equation~\eqref{eq:tail-tableaux} as
$(m-t,\nu)$ with $0\le t\le s$ and $\nu\vdash t$.  The hook-length formula
shows that $f^{(m-t,\nu)}$ is eventually a polynomial in $m$ of degree $t$
with leading coefficient $f^\nu/t!$.  Consequently $E_s(m)$ is eventually
a polynomial of degree $s$ with leading coefficient
$\sum_{\nu\vdash s}f^\nu/s!=I_s/s!$.  Thus the $j=0$ term in
Equation~\eqref{eq:near-stable-defect} supplies degree $r-1$ with the stated
leading coefficient, while every later term has smaller degree.  Finally,
direct hook-length evaluation gives
\[
 E_0(m)=1,\quad E_1(m)=m,\quad E_2(m)=(m-1)^2,\quad
 E_3(m)=\frac{m(m-2)(2m-5)}3,
\]
in the ranges used above.  Substitution proves
Equation~\eqref{eq:first-near-stable-strips}.
\end{proof}

In particular, the stability threshold $n\ge k$ is sharp for every $k\ge3$:
the immediately preceding rank has exactly one fewer object.

The stable sequence is OEIS~A170941 \cite{OEISA170941}.  Substitution of the
involution recurrence into Equation~\eqref{eq:stable-ie}, followed by
Pascal's identity, gives the useful recurrence
\begin{equation}\label{eq:stable-recurrence}
 a_{k+1}=a_k+k a_{k-1}-a_{k-2}+a_{k-3}\qquad(k\ge3),
\end{equation}
with $a_0=a_1=a_2=1$ and $a_3=2$.

For $n=3$, the standard arc representation identifies the peakless Motzkin
specialization with noncrossing involutions having no adjacent transposition;
see also the RNA arc-diagram setting of \cite{GilLopez2023}.

\section{The limiting proportion}

Let $\Pi_k$ be a uniformly random involution on $[k]$, and let
\[
 X_k=\#\{i:(i,i+1)\text{ is a cycle of }\Pi_k\}.
\]

\begin{theorem}[Poisson limit]\label{thm:poisson}
As $k\to\infty$,
\[
 X_k\ \xrightarrow{d}\ \operatorname{Poisson}(1).
\]
Consequently
\begin{equation}\label{eq:e-minus-one}
 \frac{a_k}{I_k}\longrightarrow e^{-1}.
\end{equation}
\end{theorem}

\begin{proof}
For fixed $r$, the falling factorial $(X_k)_r$ orders $r$ distinct adjacent
transpositions.  They can occur simultaneously exactly when their path edges
form a matching.  Fixing them leaves an arbitrary involution on the remaining
labels, so
\begin{equation}\label{eq:factorial-moments}
 \mathbb E[(X_k)_r]
 =r!\binom{k-r}{r}\frac{I_{k-2r}}{I_k}.
\end{equation}
The classical involution asymptotic of Moser and Wyman
\cite{MoserWyman1955} implies, for fixed $r$,
\[
 \frac{I_{k-2r}}{I_k}\sim k^{-r}.
\]
Since $r!\binom{k-r}{r}\sim k^r$, every fixed factorial moment in
Equation~\eqref{eq:factorial-moments} tends to $1$.  These are the factorial
moments of $\operatorname{Poisson}(1)$.  The first moment makes the family
tight.  Moreover, for $M>r$,
\[
 \mathbb E\!\left[(X_k)_r\mathbf1_{\{X_k>M\}}\right]
 \le \frac{\mathbb E[(X_k)_{r+1}]}{M-r},
\]
so the bounded $(r+1)$st factorial moments make $(X_k)_r$ uniformly
integrable.  Every subsequential limit therefore has all factorial moments
equal to one and hence is $\operatorname{Poisson}(1)$.  This proves the
distributional limit.  Finally, evaluating the distribution functions at
the continuity point $1/2$ gives
$\mathbb P(X_k=0)\to e^{-1}$.  The event $X_k=0$ is precisely the class
counted by $a_k$, yielding Equation~\eqref{eq:e-minus-one}.
\end{proof}

Equation~\eqref{eq:e-minus-one} proves the limiting-proportion conjecture
recorded in OEIS~A170941.  It also gives the stable ribbon asymptotic
immediately from the Moser--Wyman formula.

Combining the fixed-rank and stable results distinguishes the two asymptotic
regimes.

\begin{corollary}[Finite-rank density and noncommuting limits]
For every fixed $n\ge3$,
\begin{equation}\label{eq:fixed-rank-density}
 \frac{b_{n,k}}{u_{n,k}}\sim
 \frac{\alpha_n}{n}
 \left(\frac{\sqrt{n^2-4}}n\right)^{\beta_n-1}
 \left(\frac{\alpha_n}{n}\right)^k.
\end{equation}
In particular,
\begin{equation}\label{eq:noncommuting-limits}
 \lim_{n\to\infty}\lim_{k\to\infty}\frac{b_{n,k}}{u_{n,k}}=0,
 \qquad
 \lim_{k\to\infty}\lim_{n\to\infty}\frac{b_{n,k}}{u_{n,k}}=e^{-1}.
\end{equation}
\end{corollary}

\begin{proof}
Divide Equation~\eqref{eq:fixed-rank-asymptotic} by
Equation~\eqref{eq:regev-leading}.  Since $\alpha_n<n$ for every $n\ge3$,
the fixed-rank limit is zero.  In the reverse order, fixed $k$ enters the
stable range once $n\ge k$, so the inner limit is $a_k/I_k$; then apply
Equation~\eqref{eq:e-minus-one}.
\end{proof}

\section{Further directions}

The original all-$n$ enumeration problem is resolved by
Theorem~\ref{thm:main}; Theorems~\ref{thm:highest-weight} and
\ref{thm:fixed-rank-asymptotic}, together with
Corollary~\ref{cor:near-stable}, give the endpoint refinement, the complete
fixed-rank leading asymptotic, and every fixed-distance correction below
stability.  Two natural directions remain.  First, the stable
involution model suggests that the finite-rank
numbers should admit a direct positive arc-diagram interpretation with an
$n$-dependent bounded statistic.  The ordinary
longest-decreasing-subsequence bound is not that statistic: the candidate
count is $32$ at $(n,k)=(4,6)$, whereas Theorem~\ref{thm:main} gives
$b_{4,6}=31$.  The mixed defining/dual walk therefore records a different
finite-rank quotient.  Second,
Equation~\eqref{eq:noncommuting-limits} points to a uniform double-scaling
problem between the fixed-rank and stable regimes.  The proof here keeps $n$
fixed before $k\to\infty$ and does not supply estimates uniform when both
parameters grow.

\begingroup
\footnotesize
\medskip
\noindent\textbf{Declaration of Generative AI and AI-assisted technologies in
the writing process.}\par
During the preparation of this work, the author used an author-developed local
implementation of the Proof Engine framework for AI-assisted mathematical
research.  The implementation operates through replaceable AI-agent and
harness components; the present work used OpenAI Codex within that
implementation.  The framework is described in \emph{Proof Engine
Infrastructure}~\cite{LiProofEngine2026}.  The author reviewed and edited the
resulting material and assumes full responsibility for the definitions,
statements, proofs, citations, code, and exposition.

\bibliographystyle{plain}
\bibliography{references}
\endgroup
\end{document}
